% Generated draft: requires mathematical and source review.
\documentclass[11pt]{article}
\usepackage[margin=25mm]{geometry}
\usepackage{fontspec}
\setmainfont{Latin Modern Roman}
\usepackage{amsmath,amssymb,mathtools,unicode-math}
\setmathfont{Latin Modern Math}
\usepackage{enumitem,needspace,xurl,hyperref}
\setlength{\parindent}{0pt}
\setlength{\parskip}{5pt}
\setlength{\emergencystretch}{3em}
\newcommand{\N}{\mathbb{N}}
\newcommand{\Z}{\mathbb{Z}}
\newcommand{\Q}{\mathbb{Q}}
\newcommand{\R}{\mathbb{R}}
\newcommand{\C}{\mathbb{C}}
\newcommand{\F}{\mathbb{F}}
\newcommand{\A}{\mathbb{A}}
\newcommand{\PP}{\mathbb{P}}
\newcommand{\E}{\mathbb{E}}
\newcommand{\eps}{\varepsilon}
\newcommand{\dd}{\,\mathrm{d}}
\begin{document}
\section*{UnsolvedMath ID: 800337}
\textbf{Source problem code:} catalogue entry
\subsection{Definitions and mathematical statement}
Use the following standard coordinate convention for an Aztec diamond. For an integer \(n\geq 1\), let
\[
\mathcal A_n=\left\{Q_{i,j}=[i,i+1]\times[j,j+1] : i,j\in\Z,\ \left|i+\tfrac12\right|+\left|j+\tfrac12\right|\leq n\right\}.
\]
Here the elements of \(\mathcal A_n\) are unit square cells. Two cells share an edge when their integer indices differ by one of \((1,0),(-1,0),(0,1),(0,-1)\). A domino tiling of a collection of cells is a partition of that collection into pairs of cells sharing an edge.

Let \(p=Q_{a,b}\) and \(q=Q_{c,d}\) be cells of \(\mathcal A_n\). They are related by a knight's move if
\[
(|a-c|,|b-d|)\in\{(1,2),(2,1)\}.
\]
For such a pair define
\[
T(n;p,q)=\#\left\{\text{domino tilings of }\mathcal A_n\setminus\{p,q\}\right\}.
\]
\textbf{Problem attributed to Propp (1998), Problem 8.} Determine \(T(n;p,q)\) for the two ``close-to-central'' cells intended in the source, which are related by a knight's move.

\textbf{Unresolved qualification.} The supplied record does not specify the coordinates of the deleted cells, a mathematical meaning of ``close-to-central'', any parity dependence of their positions on \(n\), or the intended range of orders. The definitions above give a coordinate framework for expressing the requested count; they do not identify the intended pair or replace the source question by a request to count every admissible pair. A precise family of deleted-cell pairs and its allowed orders must be recovered before the target enumeration is fully determined.

The supplied preceding fragment mentions a large prime factor \(73\), describes the other prime factors as reassuringly small, and suggests that further data might explain the occurrence of \(73\). Its numerical display is truncated, so no particular tiling count or factorization can be reliably associated with this remark here. The registry labels the problem open; that label is supplied metadata, not an independently verified assessment.
\subsection{Short English statement}
Count the domino tilings of an Aztec diamond after deleting two near-central cells separated by a knight's move. The supplied excerpt does not identify the cells precisely, so the intended family of regions remains unresolved.
\subsection{Sources}
\begin{itemize}
\item Local UnsolvedMath registry, permanent ID \(800337\), source problem code \texttt{catalogue entry}, title \emph{Enumeration of Matchings --- Problem 8}; supplied record in the AMR Open Problem Lists collection. The registry attributes the problem to Propp, gives the year 1998, and labels it open.
\item Propp, \emph{Twenty Open Problems in Enumeration of Matchings} (1998), Problem 8 (\#8). Supplied source URL: \url{https://arxiv.org/abs/math/9801060}. The record reports access on 2026-10-16 and extraction by \texttt{arxiv\_latex}; these are supplied provenance statements, not access or verification performed for this draft.
\item Supplied upstream revision: \texttt{d371ee917ac9647fbb06cd6fc529deec6f3dda82}. The supplied excerpt contains a truncated numerical display and a remark about the prime factor \(73\); its connection to a specific region or count cannot be reconstructed from this record alone.
\end{itemize}
\end{document}
